\documentclass[11pt]{article}

\usepackage[T1]{fontenc}
\usepackage[utf8]{inputenc}
\usepackage{amsmath}
\usepackage{amssymb}
\usepackage{amsthm}
\usepackage{mathtools}
\usepackage[hidelinks]{hyperref}
\usepackage{doi}

\setlength{\emergencystretch}{2em}

\newtheorem{theorem}{Theorem}
\newtheorem{corollary}[theorem]{Corollary}
\newtheorem{proposition}[theorem]{Proposition}
\newtheorem{lemma}[theorem]{Lemma}
\newtheorem{definition}{Definition}
\newtheorem{remark}{Remark}

\newcommand{\Heis}{\mathrm{Heis}}
\newcommand{\Cay}{\mathrm{Cay}}
\newcommand{\ran}{\mathrm{ran}}
\newcommand{\Sent}{S_{\mathrm{ent}}}
\newcommand{\Sadm}{S_{\mathrm{ent}}^{\mathrm{adm}}}
\newcommand{\rpair}{r_{\mathrm{pair}}}
\newcommand{\Hres}{[\mathrm{H}_{\mathrm{res}}]}
\newcommand{\ue}{\hat{u}}

\title{Two Entanglement Entropies of the Conjugate Weil Pair:\\
the Admissible Pair Sector and the Full Residual Rank}
\author{J\'er\^ome Beau}
\date{June 2026}

\begin{document}

  \maketitle

  \begin{abstract}
    The companion entanglement note and the orthogonality paper E1 identify two distinct
    entanglement entropies of the conjugate Weil pair $\{c,q-c\}$, which must not be
    conflated, and which now carry different epistemic status.
    The first is the entropy of the minimal admissible pair sector: the rank-3 admissible
    projection $\Pi_q$ is spectrally atomic, equal to $\mathrm{span}\{e_0,e_{\xi_c},e_{q-\xi_c}\}$;
    removing the zero mode leaves the two-dimensional Born--Infeld parity pair
    $\{e_{\xi_c},e_{q-\xi_c}\}$, identified with the spin-$\tfrac12$ representation.
    Under the explicit, undemonstrated hypothesis (H-inv) that the proto-state on this sector is
    invariant under the diagonal $2I$-action -- a hypothesis stated and left open, not derived,
    by \emph{Q3} -- the proto-state is the singlet, its reduced state is maximally mixed, and
    $S_{\mathrm{ent}}^{\mathrm{adm}}=\log 2$ for general canonical blocks.
    This admissible pair entropy is therefore conditional on (H-inv), not unconditional.
    The second object is the full residual fibre-level rank $r_{\mathrm{pair}}(n)=R_\infty-R(n)$
    defined on the full Gram--Schmidt basis (rank $q$ at saturation, not the rank-3 admissible
    projection).
    E1 proves $S_{\mathrm{ent}}(n)=\log(r_{\mathrm{pair}}(n))$ unconditionally for matched
    single-character blocks -- this is the one unconditional result of the present note; for
    independently sampled canonical blocks it remains conditional on the residual admissibility
    indiscernibility hypothesis $[\mathrm{H}_{\mathrm{res}}]$, and on extending \emph{O17}'s
    equal-residual-dimension result beyond its own toy model to the real admissibility pipeline.
    The metaplectic dilation $\phi_\omega$ is excluded as a bridge for $[\mathrm{H}_{\mathrm{res}}]$
    because it acts between blocks, not within the residual support of one block.
  \end{abstract}

  \noindent\textbf{Keywords:}
  conjugate Weil pairs;
  spectral atomicity;
  admissible projection;
  spin-$\tfrac12$ singlet;
  diagonal invariance hypothesis;
  residual rank;
  entanglement entropy;
  epistemic stratification.

  \newpage
  \tableofcontents

  \newpage

  \section{Introduction}
    \label{sec:intro}

    The companion entanglement note~\cite{Beau2026entnote} and the orthogonality paper
    E1~\cite{Beau2026ent1} establish results about the entanglement of a closed conjugate Weil
    pair $\{c,q-c\}$ in the spectral admissibility cascade.
    A careful reading shows that two genuinely different entanglement entropies are in play,
    attached to two different projections, with different epistemic status.
    Conflating them produces an apparent tension between a $\log 2$ figure and a depth-dependent
    $\log\rpair(n)$ figure that in fact rest on different hypotheses.

    This note separates them.
    The separation is not cosmetic: the two objects live on projections of different rank,
    have different physical roles, and have different epistemic status.
    One is the entropy of the minimal admissible sector, which is two-dimensional and carries
    the singlet relevant to diagonal-invariant bipartite structure; it holds for general blocks
    but only conditionally, under the explicit hypothesis (H-inv) that \emph{Q3} states and
    leaves open.
    The other is the entropy of the full residual Gram--Schmidt rank, a finer pipeline-level
    observable; it is unconditional only for matched blocks and otherwise rests on bridge
    hypotheses of its own.

  \section{Two distinct objects}
    \label{sec:objects}

    The O-series pipeline stores two distinct objects per conjugate pair, and
    \emph{Q5a-O2}~\cite{Beau2026q5ao2} (Remark 2.1) is explicit that they must not be identified.

    \begin{itemize}
      \item The full Gram--Schmidt basis is a unitary $q\times q$ matrix whose rows span all of
      $\mathbb{C}^q$ at saturation.
      Its cumulative admitted rank is $R(n)$, growing to $R_\infty=q$ at saturation, and
      the residual fibre-level rank is
      \[
        \rpair(n) = R_\infty - R(n),
      \]
      which is non-increasing in $n$.
      This is the object of the companion note~\cite{Beau2026entnote}.
      It is \emph{not} the theoretical admissible projection.
      \item The theoretical admissible projection $\Pi_q$ is the rank-3 projection onto the first
      three Gram--Schmidt directions $\{\phi_0,\phi_1,\phi_2\}$, identified through the
      $H_{\mathrm{eff}}$-projection coordinates of \emph{Q5a-O2}~\cite{Beau2026q5ao2}, Remark 2.1.
      All spectral statements about $\Pi_q$ concern this rank-3 object.
    \end{itemize}

    We therefore distinguish the admissible pair entropy $\Sadm$, defined on $\Pi_q$
    (Section~\ref{sec:adm}), from the full residual-rank entropy $\Sent(n)$, defined on the
    full Gram--Schmidt basis (Section~\ref{sec:residual}).

  \section{The admissible pair sector: \texorpdfstring{$\Sadm=\log 2$}{S=log 2} under (H-inv)}
    \label{sec:adm}

    The admissible sector is not an arbitrary residual subspace.
    Its structure is fixed by spectral atomicity.

    \begin{theorem}[Spectral atomicity, \emph{Q5a-O2}~\cite{Beau2026q5ao2}, Proposition 4.1]
      \label{thm:atomicity}
      For each prime $q$ and each conjugate pair $(c,q-c)$, the three admissible directions
      are pure Fourier modes,
      \[
        \phi_0 = e_0,
        \qquad
        \phi_1 = e_{\xi_c},
        \qquad
        \phi_2 = e_{q-\xi_c},
      \]
      with $e_\xi(k)=q^{-1/2}e^{2\pi i\xi k/q}$ and $\xi_c\in\{1,\dots,\lfloor q/2\rfloor\}$.
      Equivalently $\ran(\Pi_q^{(c)}) = \mathrm{span}\{e_0,e_{\xi_c},e_{q-\xi_c}\}$.
    \end{theorem}

    This is an unconditional theorem of \emph{Q5a-O2}: it replaces the earlier coercivity and
    tightness hypotheses with proved statements, and is confirmed numerically to machine
    precision for $q\in\{29,61,101,151\}$.
    Its proof rests on three closed steps: the $k{=}3$ Weil fingerprints are pure Fourier
    modes (the orthogonality result of E1~\cite{Beau2026ent1}); Gram--Schmidt preserves purity;
    and the admissible rank is exactly three~\cite{Beau2026a27}, which with the Born--Infeld
    parity involution forcing $e_{\xi_c}$ and $e_{q-\xi_c}$ to enter together fixes the triple.

    The zero mode is removed by admissibility.

    \begin{lemma}[Mean-zero, \emph{Q5a-O2}~\cite{Beau2026q5ao2}, Lemma 3.1]
      \label{lem:meanzero}
      For every $f\in\ran(\Pi_q^{(c)})$, the Fourier coefficient $\hat{f}(0)=0$.
    \end{lemma}

    \begin{proof}
      The Weil representation $\rho_c$ with $c\neq 0$ is irreducible with non-trivial central
      character, so the constant function (the trivial representation) is orthogonal to
      $\ran(\Pi_q^{(c)})$, giving $\langle\mathbf{1},f\rangle = \hat{f}(0)=0$.
    \end{proof}

    After removing $e_0$, the admissible sector is the two-dimensional Born--Infeld parity pair
    \[
      \mathcal{A}_c
      := \mathrm{span}\{e_{\xi_c}, e_{q-\xi_c}\},
      \qquad
      q-\xi_c \equiv -\xi_c \pmod q,
    \]
    which we identify with the spin-$\tfrac12$ representation $V_{1/2}\cong\mathbb{C}^2$.
    The two factors of the conjugate pair $\{c,q-c\}$ each carry a copy of $\mathcal{A}_c$,
    and the bipartite proto-state lives on $\mathcal{A}_c\otimes\mathcal{A}_c$.

    \begin{quote}
      \textbf{(H-inv)} \emph{(Q3~\cite{Beau2026q3}, hypothesis, not a theorem of that paper.)}
      The proto-state $|\psi\rangle\in\mathcal{A}_c\otimes\mathcal{A}_c$ is invariant under the
      diagonal $2I$-action: for every $g\in 2I$, $(D(g)\otimes D(g))|\psi\rangle=|\psi\rangle$.
    \end{quote}

    \emph{Q3} states (H-inv) as an explicit hypothesis and proves what it buys once assumed; it
    does not derive (H-inv) from Born--Infeld indiscernibility, from admissibility, or from any
    other structural principle, and says so explicitly.
    Whether a physically admissible bipartite proto-state must satisfy (H-inv) is a separate,
    open question.
    The present note imports (H-inv) as a hypothesis on the same footing as \emph{Q3}, rather
    than deriving it.

    \begin{theorem}[Admissible pair entropy under (H-inv)]
      \label{thm:adm}
      Assume (H-inv).
      Then the proto-state on $\mathcal{A}_c\otimes\mathcal{A}_c$ is the singlet
      \[
        |\Omega_c\rangle
        = \frac{1}{\sqrt{2}}
        \bigl(|\xi_c\rangle\otimes|q-\xi_c\rangle
        - |q-\xi_c\rangle\otimes|\xi_c\rangle\bigr).
      \]
      Its reduced state is $\rho^{\mathrm{adm}}_c = \tfrac12\mathbf{1}_2$, and hence
      \[
        \Sadm = \log 2.
      \]
    \end{theorem}

    \begin{proof}
      By Theorem~\ref{thm:atomicity} and Lemma~\ref{lem:meanzero}, the admissible sector of each
      factor is $\mathcal{A}_c\cong V_{1/2}$.
      By \emph{Q3}~\cite{Beau2026q3}, Theorem 3.2, applied to $j=\tfrac12$: the invariants of the
      diagonal $2I$-action on $V_{1/2}\otimes V_{1/2}$ form a one-dimensional space (Schur's
      lemma, via $V_{1/2}\otimes V_{1/2}=V_0\oplus V_1$ and irreducibility/self-duality of
      $V_{1/2}$), spanned by the $\mathrm{SU}(2)$ singlet $|\Omega_c\rangle$.
      Under (H-inv), the proto-state lies in this space and, being normalised, equals
      $|\Omega_c\rangle$ up to phase.
      The Schmidt decomposition of $|\Omega_c\rangle$ has two terms with equal coefficients
      $1/\sqrt 2$, so $\rho^{\mathrm{adm}}_c = \tfrac12\mathbf{1}_2$ and
      $\Sadm = -\mathrm{Tr}(\tfrac12\mathbf{1}_2\log\tfrac12\mathbf{1}_2) = \log 2$.
    \end{proof}

    Theorem~\ref{thm:adm} uses only the admissible parity pair $\{e_{\xi_c},e_{q-\xi_c}\}$,
    fixed by spectral atomicity and parity for any block parameters, so it does not require a
    canonical identification of the full residual supports; that part of the argument is
    unconditional.
    What is conditional is (H-inv) itself: the admissible pair entropy is $\log 2$ for general
    canonical blocks \emph{given} (H-inv), not unconditionally.
    The present note imports both the atomicity of \emph{Q5a-O2} and the singlet identification
    of \emph{Q3} under (H-inv), rather than re-deriving either internally, and inherits the
    conditional status Q3 itself assigns to (H-inv).

  \section{The full residual rank: matched blocks and \texorpdfstring{$\Hres$}{[H res]}}
    \label{sec:residual}

    The full residual fibre-level rank $\rpair(n)=R_\infty-R(n)$ is a different object,
    defined on the full Gram--Schmidt basis rather than on $\Pi_q$.

    On matched single-character blocks, E1~\cite{Beau2026ent1} (Theorem~7) proves that the
    residual Schmidt spectrum is flat and that
    \[
      \Sent(n) = \log \rpair(n).
    \]
    The proof uses the sign-reflected Fourier supports $S_{q-c}^{\mathrm{res}}(n)
  = -S_c^{\mathrm{res}}(n)$, which canonically identify the residual modes of the two factors
    and yield a maximally entangled residual state.

    For independently sampled canonical blocks, E1 also shows that this literal mode-by-mode
    identification is not automatic.
    \emph{O17}~\cite{Beau2026a21} supplies equality of the residual dimensions of the two
    factors, but not a canonical pairing -- and only within O17's own scalar toy model, for
    an arbitrary pair of blocks: O17's own scope section states that no argument in that paper
    connects its toy model to the real admissibility pipeline, so its equal-dimension result is
    not yet shown to transfer to the conjugate-pair setting used here.
    This is carried below as a standing gap alongside $\Hres$, not folded silently into it.
    Universal admission orthogonality ($w_j=1$, E1 Theorem~4) ensures each factor has an
    orthonormal residual basis, but orthonormality and equal dimension do not by themselves
    fix the joint state or force a flat marginal.

    \begin{definition}[Residual admissibility indiscernibility $\Hres$]
      \label{def:hres}
      At BFS depth $n$, the residual fingerprint directions of a conjugate Weil factor are
      residually admissibility-indiscernible if no invariant carried by the admissible
      generators distinguishes them, equivalently if assigning unequal weights would constitute
      a premature admissible resolution of the open residual fibre.
    \end{definition}

    \begin{proposition}[Conditional flatness]
      \label{prop:hres}
      Assume the closure of the conjugate pair as a pure object, assume $\Hres$, and assume that
      \emph{O17}'s equal-residual-dimension result transfers from its own toy model to the real
      admissibility pipeline for conjugate pairs (open; not established by O17 itself).
      Then the reduced state on the residual support is
      $\rho_c^{\mathrm{res}}(n) = \rpair(n)^{-1}\mathbf{1}_{\rpair(n)}$, and
      $\Sent(n)=\log\rpair(n)$ for general canonical blocks.
    \end{proposition}

    \begin{proof}
      If unequal weights on the residual directions were admissibility-visible, they would
      define a decomposition of the open residual fibre that is stable under the admissible
      generators, i.e.\ a premature splitting, which is forbidden by Axiom A3~\cite{Beau2026m}
      and the non-factorisability lemma of the Foundation paper prior to projection locking.
      Under $\Hres$, no admissibility-visible invariant distinguishes the residual directions.
      Therefore the admissible reduced state cannot assign different weights to them without
      introducing extra structure not present in the admissibility data: the only
      admissibility-invariant reduced state is proportional to the identity.
      Together with closure (pure joint state) and the equal residual dimensions assumed to
      transfer from \emph{O17}~\cite{Beau2026a21}'s toy model, the reduced state is
      $\rho_c^{\mathrm{res}}(n)=\rpair(n)^{-1}\mathbf{1}_{\rpair(n)}$, giving
      $\Sent(n)=\log\rpair(n)$.
    \end{proof}

    \begin{remark}[The dilation does not supply $\Hres$]
      \label{rem:dilation}
      The metaplectic dilation $\phi_\omega(a,b,z)=(\omega a,\omega^{-1}b,z)$ of
      \emph{O31}~\cite{Beau2026a35} forces $\sigma_c(n)=\sigma_{\omega c}(n)$, relating the block
      of character $c$ to the block of character $\omega c$.
      It acts between blocks, not within the residual support of a single block, and so does not
      permute the residual modes of one factor.
      It cannot make the residual modes admissibility-indiscernible, and is therefore not a
      bridge for $\Hres$.
      $\Hres$ has a status parallel to the conditional hypotheses [H-color] and [H-ext] elsewhere
      in the programme.
    \end{remark}

  \section{Status}
    \label{sec:status}

    E1~\cite{Beau2026ent1} proves universal admission orthogonality: each admitted $k{=}3$ Weil
    fingerprint is a Fourier mode, and a newly admitted frequency is orthogonal to the previous
    span.
    \emph{O17}~\cite{Beau2026a21} proves equality of residual dimensions within its own toy
    model, for an arbitrary pair of blocks; extending this to the conjugate-pair setting of the
    real admissibility pipeline is open (Section~\ref{sec:residual}).

    The admissible spin-$\tfrac12$ pair sector's structure is closed unconditionally by spectral
    atomicity (\emph{Q5a-O2}~\cite{Beau2026q5ao2}, Proposition 4.1) and the mean-zero lemma: the
    zero mode is excluded and the two surviving non-zero frequencies form the parity pair
    $\{e_{\xi_c},e_{q-\xi_c}\}$, identified with $V_{1/2}$.
    What is \emph{not} unconditional is the proto-state on this sector: Theorem~\ref{thm:adm}
    identifies it as the singlet, and hence
    \[
      \Sadm = \log 2,
    \]
    only under the explicit hypothesis (H-inv) that \emph{Q3}~\cite{Beau2026q3} states and leaves
    open, not under any unconditional forcing.
    (H-inv) is not derived from Born--Infeld indiscernibility, from admissibility, or from any
    other principle in the present corpus.

    By contrast, the full residual-rank statement
    \[
      \Sent(n) = \log \rpair(n)
    \]
    is unconditional only for matched single-character blocks (E1 Theorem 7) -- the one
    unconditional result of this note.
    For independently sampled canonical blocks it remains conditional on $\Hres$, the residual
    admissibility indiscernibility hypothesis (Definition~\ref{def:hres}), and on extending
    \emph{O17}'s equal-dimension result to the real pipeline (Section~\ref{sec:residual}).
    The metaplectic dilation is excluded as a bridge for $\Hres$ (Remark~\ref{rem:dilation}).

    The separation is therefore:
    \begin{align*}
      \text{admissible pair sector, under (H-inv)}
      &\;\Longrightarrow\;
      \Sadm = \log 2,\\
      \text{residual-rank sector, matched blocks}
      &\;\Longrightarrow\;
      \Sent(n) = \log \rpair(n).
    \end{align*}
    The first holds for general canonical blocks, but only conditionally on (H-inv).
    The second holds unconditionally for matched single-character blocks, and conditionally
    otherwise, on $\Hres$ together with the O17 transfer gap.

    \newpage
    \bibliographystyle{plain}
    \bibliography{cosmochrony-bibliography}

\end{document}
